\documentclass[aspectratio=169,11pt]{beamer}
\def\courseweek{4}
% Shared style for the Introduction to Information Geometry lecture slides.
% Compile the lecture .tex files from this directory.
\providecommand{\courseweek}{1}
\usepackage[T1]{fontenc}
\usepackage[utf8]{inputenc}
\usepackage[sc]{mathpazo}
\renewcommand{\sfdefault}{lmss}
\renewcommand{\ttdefault}{lmtt}
\usepackage{amsmath,amssymb,mathtools}
\usepackage{tikz}
\usetikzlibrary{arrows.meta,calc,positioning}
\usepackage{pgfplots}
\pgfplotsset{compat=1.18}
\usepackage{booktabs}

\definecolor{courseink}{HTML}{173749}
\definecolor{coursemuted}{HTML}{526674}
\definecolor{courseteal}{HTML}{0E7287}
\definecolor{coursecopper}{HTML}{B66B43}
\definecolor{courseline}{HTML}{D4E0E3}
\definecolor{coursepale}{HTML}{EFF5F5}

\setbeamercolor{normal text}{fg=courseink,bg=white}
\setbeamercolor{structure}{fg=courseteal}
\setbeamercolor{frametitle}{fg=courseink,bg=white}
\setbeamercolor{title}{fg=courseink}
\setbeamercolor{subtitle}{fg=coursemuted}
\setbeamercolor{block title}{fg=courseink,bg=white}
\setbeamercolor{block body}{fg=courseink,bg=white}
\setbeamercolor{alerted text}{fg=coursecopper}
\setbeamerfont{title}{family=\rmfamily,series=\bfseries,size=\fontsize{32}{37}}
\setbeamerfont{frametitle}{family=\rmfamily,series=\bfseries,size=\fontsize{26}{31}}
\setbeamerfont{normal text}{size=\fontsize{18}{22}}
\setbeamerfont{subtitle}{size=\fontsize{17}{21}}
\setbeamersize{text margin left=0.78cm,text margin right=0.78cm}
\setbeamertemplate{navigation symbols}{}
\setbeamertemplate{itemize items}[circle]
\setbeamertemplate{enumerate items}[default]
\setbeamertemplate{blocks}[default]
\setlength{\parskip}{0.22em}
\setbeamertemplate{footline}{%
  \leavevmode\hbox to \paperwidth{\hspace{.78cm}%
  {\color{coursemuted}\fontsize{8.5}{10}\selectfont Introduction to Information Geometry \hspace{.32cm} / \hspace{.32cm} Week \courseweek}%
  \hfill{\color{coursemuted}\fontsize{8.5}{10}\selectfont \insertframenumber\hspace{.17cm}/\hspace{.17cm}\inserttotalframenumber}\hspace{.78cm}}%
  \vskip.17cm
}
\setbeamertemplate{frametitle}{%
  \vspace{.13cm}
  \insertframetitle\par
  \vspace{.08cm}
}
\setbeamertemplate{title page}{%
  \vspace*{1.00cm}
  {\color{courseteal}\fontsize{10}{13}\selectfont\bfseries
  INTRODUCTION TO INFORMATION GEOMETRY\par}
  \vspace{.40cm}
  {\usebeamerfont{title}\usebeamercolor[fg]{title}\inserttitle\par}
  \vspace{.22cm}
  {\usebeamerfont{subtitle}\usebeamercolor[fg]{subtitle}\insertsubtitle\par}
  \vfill
  {\color{coursemuted}\fontsize{13}{16}\selectfont Christopher Lee\par}
  \vspace{.34cm}
}

\pgfplotsset{
  courseplot/.style={
    width=11.05cm,
    height=5.15cm,
    axis line style={coursemuted!75},
    tick style={coursemuted!65},
    ticklabel style={font=\scriptsize,color=coursemuted},
    label style={font=\normalsize,color=courseink},
    grid=major,
    major grid style={courseline!70,line width=.4pt},
    line width=1.45pt,
    legend style={font=\small,draw=none,fill=none,at={(.98,.03)},anchor=south east},
    clip=true
  }
}

\newcommand{\highlightmath}[1]{{\color{courseteal}#1}}
\newcommand{\thinrule}{\par\vspace{.25cm}{\color{courseline}\rule{\linewidth}{.8pt}}\par}

% A one-dimensional state space.  x=0 and x=1 are placed at 0 and 9 units.
\newcommand{\stateinterval}[1]{%
 \begin{tikzpicture}[x=1cm,y=1cm]
   \draw[courseink,line width=1.25pt] (0,0)--(9,0);
   \foreach \pos/\lab in {0/0,9/1} {
     \fill[courseink] (\pos,0) circle(2.7pt);
     \node[below=9pt,text=coursemuted] at (\pos,0) {\(\lab\)};
   }
   #1
 \end{tikzpicture}%
}

% Paired pages are separate frames in the PDF, with matching titles.
% Supply the whole frame body for each version; the instructor may edit them
% independently without changing overlay logic.
\newcommand{\coursepair}[3]{%
  \begin{frame}{#1}#2\end{frame}%
  \begin{frame}{#1}#3\end{frame}%
}

% Equilateral representation of the categorical simplex, in barycentric
% coordinates with vertices A=(0,0), B=(1,0), C=(1/2,sqrt(3)/2).
\newcommand{\coursesimplex}[1]{%
 \begin{tikzpicture}[scale=4.9]
  \fill[coursepale] (0,0)--(1,0)--(.5,.8660254)--cycle;
  \draw[courseink,line width=.65pt] (0,0)--(1,0)--(.5,.8660254)--cycle;
  \node[below left,font=\small] at (0,0) {A};
  \node[below right,font=\small] at (1,0) {B};
  \node[above,font=\small] at (.5,.8660254) {C};
  #1
 \end{tikzpicture}%
}

\title{Curves of Distributions and Sampling}
\subtitle{Week 4 \(\cdot\) Lecture 2}
\date{}
\begin{document}
\begin{frame}[plain]\titlepage\end{frame}

\begin{frame}{A path through probability laws}
Let \(P_0\) and \(P_1\) be probability vectors. For \(0\leq t\leq1\),
\[
\highlightmath{P(t)=(1-t)P_0+tP_1}.
\]
\vspace{.42cm}
This mixture path stays in the simplex because each coordinate
is nonnegative and the coordinates sum to one.
\vspace{.42cm}
Its velocity is constant: \(P'(t)=P_1-P_0\).
\end{frame}

\begin{frame}{A straight path in a triangle}
Let \(P_0=(.6,.3,.1)\) and \(P_1=(.2,.3,.5)\).
\begin{center}\scalebox{.90}{\coursesimplex{
 \coordinate (u) at (.35,.086603);
 \coordinate (v) at (.55,.433013);
 \draw[-{Stealth[length=3mm]},courseteal,line width=1.4pt](u)--(v);
 \fill[courseteal](u)circle(1.8pt);
 \fill[coursecopper](v)circle(1.8pt);
 \node[below left,font=\scriptsize]at(u){\(P_0\)};
 \node[above right,font=\scriptsize]at(v){\(P_1\)};
}}\end{center}
\vspace{.08cm}
Every point along the segment is a valid distribution.
\end{frame}

\coursepair{The mixture coordinates}{
For \(P_0=(.6,.3,.1)\), \(P_1=(.2,.3,.5)\), compute
\vspace{.45cm}
\[
P(t)=(1-t)P_0+tP_1.
\]
\vspace{.40cm}
Write each component as a function of \(t\) and find \(P'(t)\).
}{
For \(P_0=(.6,.3,.1)\), \(P_1=(.2,.3,.5)\),
\vspace{.45cm}
\[
\highlightmath{P(t)=(.6-.4t,\ .3,\ .1+.4t)}.
\]
\vspace{.40cm}
Its velocity is \(\highlightmath{P'(t)=(-.4,0,.4)}\).
}

\begin{frame}{Differentiate expectation along a path}
Fix the outcome values \(z_i\). Along \(P(t)\),
\[
\mu(t)=\sum_i z_i p_i(t).
\]
\vspace{.45cm}
Differentiating gives
\[
\highlightmath{\mu'(t)=\sum_i z_i p_i'(t)}.
\]
\vspace{.45cm}
The value distribution changes because probability moves
among the possible outcomes.
\end{frame}

\coursepair{The rate of change of a mean}{
Let \(z=(0,1,3)\) and
\(P(t)=(.6-.4t,.3,.1+.4t)\).
\vspace{.38cm}
\[
\mu(t)=0p_A(t)+p_B(t)+3p_C(t).
\]
\vspace{.35cm}
Find \(\mu(t)\) and \(\mu'(t)\).
}{
Let \(z=(0,1,3)\) and
\(P(t)=(.6-.4t,.3,.1+.4t)\).
\vspace{.38cm}
\[
\mu(t)=.3+3(.1+.4t)=\highlightmath{.6+1.2t}.
\]
\vspace{.35cm}
Thus \(\highlightmath{\mu'(t)=1.2}\).
}

\begin{frame}{What a sample gives us}
Suppose \(n\) independent observations produce counts
\(n_A,n_B,n_C\) with total \(n\).
\vspace{.42cm}
The empirical law is
\[
\highlightmath{\widehat P=\left(\frac{n_A}{n},\frac{n_B}{n},\frac{n_C}{n}\right)}.
\]
\vspace{.40cm}
It is itself a point of the simplex. It varies from sample to sample.
\end{frame}

\coursepair{An empirical distribution}{
A sample of 20 observations records
\[
n_A=11,\quad n_B=6,\quad n_C=3.
\]
\vspace{.42cm}
Find \(\widehat P\). Compare it with the proposed model
\(P=(.6,.3,.1)\).
}{
A sample of 20 observations records
\[
n_A=11,\quad n_B=6,\quad n_C=3.
\]
\vspace{.42cm}
\[
\highlightmath{\widehat P=(.55,.30,.15)}.
\]
\vspace{.3cm}
The empirical and model distributions are nearby but unequal.
}

\begin{frame}{The expected empirical frequencies}
For independent draws from a fixed categorical law \(P\),
\[
E[\widehat p_i]=p_i.
\]
\vspace{.45cm}
Their variances are
\[
\highlightmath{\operatorname{Var}(\widehat p_i)=\frac{p_i(1-p_i)}{n}}.
\]
\vspace{.35cm}
Larger samples typically produce smaller frequency fluctuations.
\end{frame}

\begin{frame}{Sampling coordinates are linked}
Because \(\widehat p_A+\widehat p_B+\widehat p_C=1\),
an increase in one empirical frequency must be offset elsewhere.
\vspace{.45cm}
For different categories \(i\ne j\),
\[
\highlightmath{\operatorname{Cov}(\widehat p_i,\widehat p_j)
=-\frac{p_ip_j}{n}}.
\]
\vspace{.40cm}
The normalization constraint appears in sampling uncertainty
as well as in deterministic dynamics.
\end{frame}

\begin{frame}{A comparison of two probability states}
Suppose the outcomes are \(z=(0,1,3)\).
\vspace{.35cm}
\[
P=(.5,.3,.2),\qquad Q=(.3,.3,.4).
\]
\vspace{.4cm}
The change from \(P\) to \(Q\) transfers \(0.2\) of the total mass
from the outcome \(0\) to the outcome \(3\).
\vspace{.34cm}
Its expectation therefore increases by \(0.2\cdot3=0.6\).
\end{frame}

\begin{frame}{Preparing to measure information}
We can now describe a distribution, draw samples from it, and follow
smooth curves in its state space.
\vspace{.53cm}
Next, we ask two further questions:
\vspace{.25cm}
\begin{itemize}\setlength{\itemsep}{.42cm}
\item How uncertain is an observation drawn from a given distribution?
\item How distinguishable are two nearby distributions?
\end{itemize}
\vspace{.32cm}
The first leads to entropy; the second leads toward information geometry.
\end{frame}
\end{document}
